% figA9_graph_robust.tex — Figure A9: T02 graph robustness generalisation
% % figA9_graph_robust.tex — Figure A9: T02 graph robustness generalisation
% % figA9_graph_robust.tex — Figure A9: T02 graph robustness generalisation
% % figA9_graph_robust.tex — Figure A9: T02 graph robustness generalisation
% \input{figures/figA9_graph_robust}

\begin{figure*}[t]
\centering
\begin{subfigure}[b]{0.50\textwidth}
\begin{tikzpicture}
\begin{axis}[
  pmh bar,
  symbolic x coords={Clean,{Edge drop\\30\%},{Feature drop\\30\%},{Worst case}},
  xtick=data,
  xticklabel style={font=\footnotesize, align=center},
  ylabel={AUC / Accuracy (\%)},
  ymin=0, ymax=105,
  legend style={at={(0.5,1.02)}, anchor=south, legend columns=3,
                draw=none, fill=none, font=\footnotesize},
  width=0.92\linewidth, height=5.5cm,
  bar width=5pt,
]
\addplot[fill=colB0, draw=none]
  coordinates {(Clean,77.68)({Edge drop\\30\%},60.18)
               ({Feature drop\\30\%},33.04)({Worst case},31.61)};
\addlegendentry{B0}
\addplot[fill=colVAT, draw=none]
  coordinates {(Clean,75.00)({Edge drop\\30\%},53.75)
               ({Feature drop\\30\%},31.96)({Worst case},31.61)};
\addlegendentry{VAT}
\addplot[fill=colPMH, draw=none]
  coordinates {(Clean,79.46)({Edge drop\\30\%},76.43)
               ({Feature drop\\30\%},71.61)({Worst case},63.75)};
\addlegendentry{E1 (PMH)}
\node[font=\tiny\bfseries, color=colPMH, align=center]
  at (rel axis cs:0.65,0.86) {Feature drop:\\PMH vs B0 $+38.6$\,pp};
\end{axis}
\end{tikzpicture}
\caption{Robustness to unseen graph perturbation types.}
\end{subfigure}
\hfill
\begin{subfigure}[b]{0.46\textwidth}
\begin{tikzpicture}
\begin{axis}[
  pmh bar,
  symbolic x coords={Consistency,{AUC curve}},
  xtick=data,
  xticklabels={Prediction\\consistency, AUC under\\noise curve},
  xticklabel style={font=\footnotesize, align=center},
  ylabel={\%},
  ymin=0, ymax=100,
  legend style={at={(0.5,1.02)}, anchor=south, legend columns=3,
                draw=none, fill=none, font=\footnotesize},
  width=0.92\linewidth, height=5.5cm,
  bar width=7pt,
]
\addplot[fill=colB0,  draw=none]
  coordinates {(Consistency,66.07)({AUC curve},68.30)};
\addlegendentry{B0}
\addplot[fill=colVAT, draw=none]
  coordinates {(Consistency,58.93)({AUC curve},62.05)};
\addlegendentry{VAT}
\addplot[fill=colPMH, draw=none]
  coordinates {(Consistency,81.25)({AUC curve},78.62)};
\addlegendentry{E1 (PMH)}
\end{axis}
\end{tikzpicture}
\caption{Consistency and AUC under noise.}
\end{subfigure}
\caption{\textbf{Task~02 (PROTEINS): robustness generalisation to unseen perturbation types.}
\emph{(a)} PMH was trained with Gaussian node-feature noise yet generalises dramatically to
edge removal (E1 76.4\% vs.\ B0 60.2\%, $+16.2$\,pp) and feature dropout
(E1 71.6\% vs.\ B0 33.0\%, $+38.6$\,pp) --- perturbation types never seen during training.
Worst-case accuracy: E1 63.75\% vs.\ B0 31.61\%.
\emph{(b)} PMH achieves 81.3\% prediction consistency and 78.6\% AUC under the noise curve,
vs.\ B0 66.1\% and 68.3\%. Robustness generalisation is a by-product of global Frobenius
regularisation, confirming the geometric repair interpretation of PMH.}
\label{fig:A9}
\end{figure*}


\begin{figure*}[t]
\centering
\begin{subfigure}[b]{0.50\textwidth}
\begin{tikzpicture}
\begin{axis}[
  pmh bar,
  symbolic x coords={Clean,{Edge drop\\30\%},{Feature drop\\30\%},{Worst case}},
  xtick=data,
  xticklabel style={font=\footnotesize, align=center},
  ylabel={AUC / Accuracy (\%)},
  ymin=0, ymax=105,
  legend style={at={(0.5,1.02)}, anchor=south, legend columns=3,
                draw=none, fill=none, font=\footnotesize},
  width=0.92\linewidth, height=5.5cm,
  bar width=5pt,
]
\addplot[fill=colB0, draw=none]
  coordinates {(Clean,77.68)({Edge drop\\30\%},60.18)
               ({Feature drop\\30\%},33.04)({Worst case},31.61)};
\addlegendentry{B0}
\addplot[fill=colVAT, draw=none]
  coordinates {(Clean,75.00)({Edge drop\\30\%},53.75)
               ({Feature drop\\30\%},31.96)({Worst case},31.61)};
\addlegendentry{VAT}
\addplot[fill=colPMH, draw=none]
  coordinates {(Clean,79.46)({Edge drop\\30\%},76.43)
               ({Feature drop\\30\%},71.61)({Worst case},63.75)};
\addlegendentry{E1 (PMH)}
\node[font=\tiny\bfseries, color=colPMH, align=center]
  at (rel axis cs:0.65,0.86) {Feature drop:\\PMH vs B0 $+38.6$\,pp};
\end{axis}
\end{tikzpicture}
\caption{Robustness to unseen graph perturbation types.}
\end{subfigure}
\hfill
\begin{subfigure}[b]{0.46\textwidth}
\begin{tikzpicture}
\begin{axis}[
  pmh bar,
  symbolic x coords={Consistency,{AUC curve}},
  xtick=data,
  xticklabels={Prediction\\consistency, AUC under\\noise curve},
  xticklabel style={font=\footnotesize, align=center},
  ylabel={\%},
  ymin=0, ymax=100,
  legend style={at={(0.5,1.02)}, anchor=south, legend columns=3,
                draw=none, fill=none, font=\footnotesize},
  width=0.92\linewidth, height=5.5cm,
  bar width=7pt,
]
\addplot[fill=colB0,  draw=none]
  coordinates {(Consistency,66.07)({AUC curve},68.30)};
\addlegendentry{B0}
\addplot[fill=colVAT, draw=none]
  coordinates {(Consistency,58.93)({AUC curve},62.05)};
\addlegendentry{VAT}
\addplot[fill=colPMH, draw=none]
  coordinates {(Consistency,81.25)({AUC curve},78.62)};
\addlegendentry{E1 (PMH)}
\end{axis}
\end{tikzpicture}
\caption{Consistency and AUC under noise.}
\end{subfigure}
\caption{\textbf{Task~02 (PROTEINS): robustness generalisation to unseen perturbation types.}
\emph{(a)} PMH was trained with Gaussian node-feature noise yet generalises dramatically to
edge removal (E1 76.4\% vs.\ B0 60.2\%, $+16.2$\,pp) and feature dropout
(E1 71.6\% vs.\ B0 33.0\%, $+38.6$\,pp) --- perturbation types never seen during training.
Worst-case accuracy: E1 63.75\% vs.\ B0 31.61\%.
\emph{(b)} PMH achieves 81.3\% prediction consistency and 78.6\% AUC under the noise curve,
vs.\ B0 66.1\% and 68.3\%. Robustness generalisation is a by-product of global Frobenius
regularisation, confirming the geometric repair interpretation of PMH.}
\label{fig:A9}
\end{figure*}


\begin{figure*}[t]
\centering
\begin{subfigure}[b]{0.50\textwidth}
\begin{tikzpicture}
\begin{axis}[
  pmh bar,
  symbolic x coords={Clean,{Edge drop\\30\%},{Feature drop\\30\%},{Worst case}},
  xtick=data,
  xticklabel style={font=\footnotesize, align=center},
  ylabel={AUC / Accuracy (\%)},
  ymin=0, ymax=105,
  legend style={at={(0.5,1.02)}, anchor=south, legend columns=3,
                draw=none, fill=none, font=\footnotesize},
  width=0.92\linewidth, height=5.5cm,
  bar width=5pt,
]
\addplot[fill=colB0, draw=none]
  coordinates {(Clean,77.68)({Edge drop\\30\%},60.18)
               ({Feature drop\\30\%},33.04)({Worst case},31.61)};
\addlegendentry{B0}
\addplot[fill=colVAT, draw=none]
  coordinates {(Clean,75.00)({Edge drop\\30\%},53.75)
               ({Feature drop\\30\%},31.96)({Worst case},31.61)};
\addlegendentry{VAT}
\addplot[fill=colPMH, draw=none]
  coordinates {(Clean,79.46)({Edge drop\\30\%},76.43)
               ({Feature drop\\30\%},71.61)({Worst case},63.75)};
\addlegendentry{E1 (PMH)}
\node[font=\tiny\bfseries, color=colPMH, align=center]
  at (rel axis cs:0.65,0.86) {Feature drop:\\PMH vs B0 $+38.6$\,pp};
\end{axis}
\end{tikzpicture}
\caption{Robustness to unseen graph perturbation types.}
\end{subfigure}
\hfill
\begin{subfigure}[b]{0.46\textwidth}
\begin{tikzpicture}
\begin{axis}[
  pmh bar,
  symbolic x coords={Consistency,{AUC curve}},
  xtick=data,
  xticklabels={Prediction\\consistency, AUC under\\noise curve},
  xticklabel style={font=\footnotesize, align=center},
  ylabel={\%},
  ymin=0, ymax=100,
  legend style={at={(0.5,1.02)}, anchor=south, legend columns=3,
                draw=none, fill=none, font=\footnotesize},
  width=0.92\linewidth, height=5.5cm,
  bar width=7pt,
]
\addplot[fill=colB0,  draw=none]
  coordinates {(Consistency,66.07)({AUC curve},68.30)};
\addlegendentry{B0}
\addplot[fill=colVAT, draw=none]
  coordinates {(Consistency,58.93)({AUC curve},62.05)};
\addlegendentry{VAT}
\addplot[fill=colPMH, draw=none]
  coordinates {(Consistency,81.25)({AUC curve},78.62)};
\addlegendentry{E1 (PMH)}
\end{axis}
\end{tikzpicture}
\caption{Consistency and AUC under noise.}
\end{subfigure}
\caption{\textbf{Task~02 (PROTEINS): robustness generalisation to unseen perturbation types.}
\emph{(a)} PMH was trained with Gaussian node-feature noise yet generalises dramatically to
edge removal (E1 76.4\% vs.\ B0 60.2\%, $+16.2$\,pp) and feature dropout
(E1 71.6\% vs.\ B0 33.0\%, $+38.6$\,pp) --- perturbation types never seen during training.
Worst-case accuracy: E1 63.75\% vs.\ B0 31.61\%.
\emph{(b)} PMH achieves 81.3\% prediction consistency and 78.6\% AUC under the noise curve,
vs.\ B0 66.1\% and 68.3\%. Robustness generalisation is a by-product of global Frobenius
regularisation, confirming the geometric repair interpretation of PMH.}
\label{fig:A9}
\end{figure*}


\begin{figure*}[t]
\centering
\begin{subfigure}[b]{0.50\textwidth}
\begin{tikzpicture}
\begin{axis}[
  pmh bar,
  symbolic x coords={Clean,{Edge drop\\30\%},{Feature drop\\30\%},{Worst case}},
  xtick=data,
  xticklabel style={font=\footnotesize, align=center},
  ylabel={AUC / Accuracy (\%)},
  ymin=0, ymax=105,
  legend style={at={(0.5,1.02)}, anchor=south, legend columns=3,
                draw=none, fill=none, font=\footnotesize},
  width=0.92\linewidth, height=5.5cm,
  bar width=5pt,
]
\addplot[fill=colB0, draw=none]
  coordinates {(Clean,77.68)({Edge drop\\30\%},60.18)
               ({Feature drop\\30\%},33.04)({Worst case},31.61)};
\addlegendentry{B0}
\addplot[fill=colVAT, draw=none]
  coordinates {(Clean,75.00)({Edge drop\\30\%},53.75)
               ({Feature drop\\30\%},31.96)({Worst case},31.61)};
\addlegendentry{VAT}
\addplot[fill=colPMH, draw=none]
  coordinates {(Clean,79.46)({Edge drop\\30\%},76.43)
               ({Feature drop\\30\%},71.61)({Worst case},63.75)};
\addlegendentry{E1 (PMH)}
\node[font=\tiny\bfseries, color=colPMH, align=center]
  at (rel axis cs:0.65,0.86) {Feature drop:\\PMH vs B0 $+38.6$\,pp};
\end{axis}
\end{tikzpicture}
\caption{Robustness to unseen graph perturbation types.}
\end{subfigure}
\hfill
\begin{subfigure}[b]{0.46\textwidth}
\begin{tikzpicture}
\begin{axis}[
  pmh bar,
  symbolic x coords={Consistency,{AUC curve}},
  xtick=data,
  xticklabels={Prediction\\consistency, AUC under\\noise curve},
  xticklabel style={font=\footnotesize, align=center},
  ylabel={\%},
  ymin=0, ymax=100,
  legend style={at={(0.5,1.02)}, anchor=south, legend columns=3,
                draw=none, fill=none, font=\footnotesize},
  width=0.92\linewidth, height=5.5cm,
  bar width=7pt,
]
\addplot[fill=colB0,  draw=none]
  coordinates {(Consistency,66.07)({AUC curve},68.30)};
\addlegendentry{B0}
\addplot[fill=colVAT, draw=none]
  coordinates {(Consistency,58.93)({AUC curve},62.05)};
\addlegendentry{VAT}
\addplot[fill=colPMH, draw=none]
  coordinates {(Consistency,81.25)({AUC curve},78.62)};
\addlegendentry{E1 (PMH)}
\end{axis}
\end{tikzpicture}
\caption{Consistency and AUC under noise.}
\end{subfigure}
\caption{\textbf{Task~02 (PROTEINS): robustness generalisation to unseen perturbation types.}
\emph{(a)} PMH was trained with Gaussian node-feature noise yet generalises dramatically to
edge removal (E1 76.4\% vs.\ B0 60.2\%, $+16.2$\,pp) and feature dropout
(E1 71.6\% vs.\ B0 33.0\%, $+38.6$\,pp) --- perturbation types never seen during training.
Worst-case accuracy: E1 63.75\% vs.\ B0 31.61\%.
\emph{(b)} PMH achieves 81.3\% prediction consistency and 78.6\% AUC under the noise curve,
vs.\ B0 66.1\% and 68.3\%. Robustness generalisation is a by-product of global Frobenius
regularisation, confirming the geometric repair interpretation of PMH.}
\label{fig:A9}
\end{figure*}
